\section{Conclusion}\label{sec:conclusion}

This paper examines whether team gender composition is associated with systematic differences in how economists write. Using 9,117 abstracts and an LLM-based 16-item rubric, we first show that the LLM readability score aligns with established readability measures used by \textcite{hengel2022publishing}. We then document broad mean differences in writing style across team gender composition and show, in controlled regressions, that several of these differences remain after accounting for journal, year, and a large set of article and author characteristics.

The most robust differences are not in emotionality or generalized lack of confidence. Instead, they are concentrated in clarity and presentation: female-majority teams write more readable and direct abstracts, use less jargon, generally provide more evidentiary support, and make weaker novelty claims. Other dimensions, such as assertiveness, qualifiers, and emotional valence, show small or unstable differences.

Extending the analysis to peer review, we find that LLM readability does not exhibit the differential gender effect \textcite{hengel2022publishing} document for their readability measures: peer review improves LLM scores on average, but similarly for male- and female-authored papers. Tests of cumulative exposure likewise fail to replicate, at the aggregate level, the ``preemptive adaptation'' pattern in which experienced female authors submit cleaner drafts rather than relying on review to close the gap. Patterns resembling that mechanism do, however, reappear in specific dimensions of authorial stance: acknowledgement of limitations and modal verb strength show apparent preemptive softening, while novelty-claim strength instead shows peer review consistently dampening female authors' novelty framing regardless of experience, a pattern that could be read as persistent reviewer resistance rather than rational adaptation. We treat these dimension-specific readings as suggestive rather than conclusive: the same explanations we invoke for the aggregate null (LLM evaluation noise, limited rubric granularity, and the LLM capturing a dimension of writing distinct from review-sensitive readability) apply here as well. With that caveat, the findings raise the possibility that Hengel's mechanism is not entirely absent from LLM-evaluated writing but concentrated in narrower dimensions of stance and hedging rather than overall readability.

The paper does not identify whether these differences reflect selection, training, audience expectations, or differential review pressure, and the LLM-based measures remain subject to measurement error. Even so, the results highlight writing style as a meaningful and understudied margin along which research may be evaluated, circulated, and absorbed within the discipline.